\begin{figure}[t]
\centering
\begin{tikzpicture}[
  font=\scriptsize,
  block/.style={rounded corners=2pt, draw=black!70, fill=blue!6, minimum width=1.05cm, minimum height=0.45cm, align=center},
  gnode/.style={circle, draw=black!70, minimum size=0.34cm, inner sep=0.5pt, align=center},
  task/.style={rounded corners=2pt, draw=black!80, fill=black!8, minimum width=0.95cm, minimum height=0.4cm, align=center},
  art/.style={gnode, fill=green!12},
  call/.style={gnode, fill=orange!16},
  st/.style={gnode, fill=purple!12},
  fail/.style={gnode, fill=red!16},
  intent/.style={rounded corners=2pt, draw=black!70, fill=yellow!18, minimum width=0.9cm, minimum height=0.38cm, align=center},
  rel/.style={->, >=Stealth, draw=black!55, thin},
  seq/.style={->, >=Stealth, draw=black!75},
]

% ============ Left panel: PTC linear history ============
\node[align=center, font=\scriptsize\bfseries] at (1.7,2.55) {PTC: Linear Execution History};

\node[block] (b1) at (0.55,1.7) {Block $P_1$};
\node[block] (b2) at (0.55,0.9) {Block $P_2$};
\node[block] (b3) at (0.55,0.1) {Block $P_3$};
\node[draw=black!30, fill=black!3, rounded corners=1pt, minimum width=1.35cm, minimum height=0.4cm, align=center] (o1) at (2.35,1.7) {\tiny stdout+err};
\node[draw=black!30, fill=black!3, rounded corners=1pt, minimum width=1.35cm, minimum height=0.4cm, align=center] (o2) at (2.35,0.9) {\tiny stdout+err};
\node[draw=black!30, fill=black!3, rounded corners=1pt, minimum width=1.35cm, minimum height=0.4cm, align=center] (o3) at (2.35,0.1) {\tiny stdout+err};
\draw[seq] (b1)--(o1); \draw[seq] (b2)--(o2); \draw[seq] (b3)--(o3);
\draw[seq] (o1) to[out=-90,in=90] (b2);
\draw[seq] (o2) to[out=-90,in=90] (b3);
\node[align=center, text=red!60!black, font=\tiny] at (1.7,-0.6) {cross-block relations flattened into text,\\ progress and dependencies implicit};

% divider
\draw[black!25, dashed] (3.5,-0.85) -- (3.5,2.75);

% ============ Right panel: GraphPTC ============
\node[align=center, font=\scriptsize\bfseries] at (8.05,2.55) {\method: Dynamic Execution Graph};

\node[task] (tk) at (8.05,1.95) {task};
\node[intent] (in) at (5.5,1.95) {intent\\ \tiny (a,z,e)};

\node[block] (gb1) at (5.5,1.15) {Block $P_1$};
\node[block] (gb2) at (5.5,0.15) {Block $P_2$};

\node[call] (c1) at (6.7,1.5) {\tiny r};
\node[call] (c2) at (6.7,0.85) {\tiny w};
\node[art]  (a1) at (7.7,1.5) {\tiny $v_1$};
\node[st]   (s1) at (7.7,0.85) {\tiny st};
\node[call] (c3) at (6.7,0.15) {\tiny r};
\node[art]  (a2) at (7.7,0.15) {\tiny $v_2$};
\node[fail] (f1) at (6.0,-0.55) {\tiny F};

% intent binding
\draw[rel] (in)--(gb1) node[midway,left,font=\tiny]{impl.};
\draw[rel, dashed] (in) to[out=0,in=180] (tk);

% block1 internals
\draw[rel] (gb1)--(c1) node[midway,above,font=\tiny]{exec};
\draw[rel] (gb1)--(c2);
\draw[rel] (c1)--(a1) node[midway,above,font=\tiny]{prod};
\draw[rel] (c2)--(s1) node[midway,above,font=\tiny]{mut};

% block2 internals + cross-block
\draw[rel] (gb2)--(c3);
\draw[rel] (c3)--(a2);
\draw[rel] (a1) to[out=-90,in=90] (c3) node[midway,right,font=\tiny]{consumes};
\draw[rel, dotted] (a2) to[out=200,in=-20] (a1) node[midway,below,font=\tiny]{equiv.};
\draw[rel] (gb2)--(f1) node[midway,left,font=\tiny]{fails\_at};

% blocks belong to task
\draw[rel, black!35] (tk) to[out=-120,in=60] (gb1);
\draw[rel, black!35] (tk) to[out=-110,in=70] (gb2);

% PATCH loop back
\draw[->, >=Stealth, red!65!black, thick] (f1) to[out=180,in=-140] node[midway,below,font=\tiny,text=red!65!black]{PATCH} (gb2.west);

\node[align=center, text=green!45!black, font=\tiny] at (8.05,-0.7) {typed in-block effects + multi-relation\\ cross-block graph, addressable and revisitable};

\end{tikzpicture}
\vspace{-2mm}
\caption{PTC records execution as a linear stream of code and output where cross-block dependencies and progress are lost, whereas \method maintains a dynamic execution graph that models the internal effects of each block (tool calls with read $r$ or write $w$ effects, artifacts $v$, states st, failures F) and cross-block dependencies (consumes, equivalence, supersession), over which a declared intent selects continue, patch, replan, or answer.}
\label{fig:overview}
\vspace{-3mm}
\end{figure}
